% TOP_PROBLEM: {"id": 1184, "title": "The abc Conjecture", "category_id": 1, "category": {"id": 1, "name": "number_theory", "display_name": "Number Theory", "description": "Properties of integers, prime numbers, Diophantine equations.", "slug": "number-theory", "order_index": 1, "created_at": "2026-07-31T15:26:25.670Z"}, "status": "open"}
% FIRST_POSTED: null
% ENTRY_KIND: statement_only
% Imported from: batch04.tex; UnsolvedMath ID 1184
\documentclass[11pt]{article}
\usepackage[margin=25mm]{geometry}
\usepackage{fontspec}
\setmainfont{Latin Modern Roman}
\usepackage{amsmath,amssymb,mathtools}
\usepackage{unicode-math}
\setmathfont{Latin Modern Math}
\usepackage{xurl,hyperref}
\hypersetup{colorlinks=true,urlcolor=blue}
\setcounter{secnumdepth}{0}
\setlength{\parindent}{0pt}
\setlength{\parskip}{5pt}
\setlength{\emergencystretch}{3em}
\newcommand{\sref}[2]{\hyperlink{source:#1:#2}{[S#2]}}
\newcommand{\cataloguescope}[1]{\paragraph{Recorded target.}#1}
\begin{document}
% UnsolvedMath ID 1184; source catalogue code NT-030
\setcounter{section}{1183}
\section{The abc Conjecture}\label{1184}
{\small\sffamily UnsolvedMath ID 1184 · Number Theory}
\subsection{Definitions and mathematical statement}

\paragraph{Shared notation.}$\mathbb N=\{1,2,\ldots\}$, and $\mathbb Z,\mathbb Q,\mathbb R,\mathbb C$ denote the integers, rationals, reals and complex numbers. Any other conventions are specified locally.


For a positive integer \(n\), define its radical by
\[
 \operatorname{rad}(n)=\prod_{\substack{p\mid n\\p\text{ prime}}}p,
 \qquad \operatorname{rad}(1)=1.
\]
The product counts each distinct prime divisor only once.
If \(a+b=c\) and \(\gcd(a,b)=1\), the three positive integers
\(a,b,c\) are pairwise coprime.
\paragraph{Statement recorded in the supplied source.}
For every real number \(\varepsilon>0\), is there a constant
\(C_\varepsilon>0\) such that every triple of positive integers
\((a,b,c)\) with \(a+b=c\) and \(\gcd(a,b)=1\) satisfies
\[
 c\le C_\varepsilon\,\operatorname{rad}(abc)^{\,1+\varepsilon}?
\]
Equivalently, for every \(\varepsilon>0\), are there only finitely many
such triples with \(c>\operatorname{rad}(abc)^{1+\varepsilon}\)? \sref{1184}{2}
\subsection{Short English statement}
Apart from finitely many exceptions for each fixed positive epsilon, is c at most the radical of abc raised to the power one plus epsilon, when a and b are coprime positive integers with a+b=c?
\subsection{Sources}
\begin{description}
\item[\hypertarget{source:1184:1}{S1}] ULAM.AI and the UnsolvedMath Contributors, \emph{UnsolvedMath: A Curated Collection of Open Mathematics Problems}, record 1184 (NT-030). CC BY 4.0. \url{https://huggingface.co/datasets/ulamai/UnsolvedMath}.
\item[\hypertarget{source:1184:2}{S2}] Jared Duker Lichtman, \emph{The abc conjecture is true almost always}, arXiv:2505.13991 (2025), Section 1, the quantified conjecture preceding (1.1). \url{https://arxiv.org/abs/2505.13991}.
\end{description}
\label{end:1184}
\end{document}
